
\section{计数伪全纯截面}\label{sec:8-6}
\zhindex{jishuweiquanchunjiemian}{计数伪全纯截面}

现在来考察与辛纤维丛相联系的\BookRecordIndex{T0143}格罗莫夫--威滕不变量。设
$(\widetilde M,\pi,\widetilde\omega)$ 是以紧黎曼曲面
$(\Sigma,j_\Sigma,\mathrm{dvol}_\Sigma)$ 为底的\BookRecordIndex{T0200}局部哈密顿纤维化，其中 $\Sigma$ 的亏格为
$g$，各纤维 $(M_z,\omega_z)$ 都辛同胚于某个紧辛 $2n$ 维流形 $(M,\omega)$。以下始终
假设各纤维满足强\BookRecordIndex{T0058}半正性条件~\eqref{eq:8-5-1}；这个条件保证，与 $M$ 上一般的二参数
$\omega$-\BookRecordIndex{T0372}驯服\BookRecordIndex{T0088}殆复结构族相联系的所有伪全纯球面都具有非负 Chern 数。选取
$\widetilde A\in H_2(\widetilde M;\mathbb Z)$，使
$\pi_*\widetilde A=[\Sigma]$；令 $k$ 为非负整数，并令
$(J,H)\in\mathcal{JH}_{\mathrm{reg}}(\widetilde M,\pi,\widetilde\omega)$ 是
定义~\ref{def:8-5-2} 意义下的正则偶对。于是，由定理~\ref{thm:8-5-1}，\BookRecordIndex{T0283}求值映射
\[
  \widetilde{\operatorname{ev}}_{J,H}:
  \mathcal M_{\Sigma,k}(\widetilde A;J,H)\longrightarrow\widetilde M^k
\]
是维数
\[
  \mu(\widetilde A,g,k)
  =n(2-2g)+2c_1^{\mathrm{Vert}}(\widetilde A)+2k
\]
的\BookRecordIndex{T0337}伪循环。由此可如下定义伪全纯截面的\BookRecordIndex{T0143}格罗莫夫--威滕不变量。

\begin{theorem}\label{thm:8-6-1}
设 $(M,\omega)$ 是满足条件~\eqref{eq:8-5-1} 的闭辛 $2n$ 维流形，而
$(\widetilde M,\pi,\widetilde\omega)$ 是以紧黎曼曲面
$(\Sigma,j_\Sigma,\mathrm{dvol}_\Sigma)$ 为底、各纤维辛同胚于 $(M,\omega)$ 的\BookRecordIndex{T0200}局部哈密顿
纤维化。设 $\widetilde A\in H_2(\widetilde M;\mathbb Z)$ 满足
$\pi_*\widetilde A=[\Sigma]$，$k$ 是非负整数，并且
$(J,H)\in\mathcal{JH}_{\mathrm{reg}}(\widetilde M,\pi,\widetilde\omega)$。则由
\begin{equation}\label{eq:8-6-1}
  \operatorname{GW}^{\widetilde M}_{\widetilde A,k}(\widetilde a)
  :=\widetilde f\mathbin{\boldsymbol\cdot}
    \widetilde{\operatorname{ev}}_{J,H}
\end{equation}
定义的群同态
\[
  \operatorname{GW}^{\widetilde M}_{\widetilde A,k}:
  H^*(\widetilde M^k)\longrightarrow\mathbb Z
\]
与正则偶对
$(J,H)\in\mathcal{JH}_{\mathrm{reg}}(\widetilde M,\pi,\widetilde\omega)$ 以及定义中
所用的\BookRecordIndex{T0337}伪循环
$\widetilde f:U\to\widetilde M^k$ 无关；这里 $\widetilde f$ 与
$\widetilde a$ Poincaré 对偶。
\end{theorem}

\begin{proof}
由定理~\ref{thm:8-5-1}，$\widetilde{\operatorname{ev}}_{J,H}$ 的配边类与 $J,H$
无关。因此，由引理~\ref{lem:6-5-5}，交数
$\widetilde f\mathbin{\boldsymbol\cdot}\widetilde{\operatorname{ev}}_{J,H}$ 与
$J,H$ 无关。再由引理~\ref{lem:6-5-7}，它只依赖于\BookRecordIndex{T0337}伪循环 $\widetilde f$ 所表示的
同调类。这就证明了定理。
\end{proof}

式~\eqref{eq:8-6-1} 定义的不变量称为 $\widetilde M$ 中类 $\widetilde A$ 的
\emph{$k$ 点伪全纯截面的\BookRecordIndex{T0143}格罗莫夫--威滕不变量}。对
$\widetilde a_1,\ldots,\widetilde a_k\in H^*(\widetilde M)$，简记
\[
  \operatorname{GW}^{\widetilde M}_{\widetilde A,k}
  (\widetilde a_1,\ldots,\widetilde a_k)
  :=\operatorname{GW}^{\widetilde M}_{\widetilde A,k}
  \bigl(\pi_1^*\widetilde a_1\smile\cdots\smile
  \pi_k^*\widetilde a_k\bigr),
\]
其中 $\pi_i:\widetilde M^k\to\widetilde M$ 表示到第 $i$ 个因子的投影。这个不变量在
几何上可作如下解释。固定一般位置的 $k$ 个循环
$\widetilde X_i\subset\widetilde M$，使它们分别与上同调类 $\widetilde a_i$ Poincaré
对偶。则该不变量带符号计数元组
$(\widetilde u,z_1,\ldots,z_k)$ 的个数，其中
$\widetilde u:\Sigma\to\widetilde M$ 是 $\widetilde J_H$-全纯截面，各 $z_i$ 是
$\Sigma$ 上两两不同的点，并且 $\widetilde u(z_i)\in\widetilde X_i$。

如下文将见，$\operatorname{GW}^{\widetilde M}_{\widetilde A,k}$ 所计数的曲线适配于
\BookRecordIndex{T0279}纤维化这一事实，赋予这些不变量额外的结构。用同调而非上同调来说明这一点最为方便。
因此，以下大部分讨论采用注~\ref{rem:7-1-12} 中的同调不变量：给定
$\widetilde M$ 中的同调类
$\widetilde\alpha_1,\ldots,\widetilde\alpha_k$，令
$\widetilde a_i:=\operatorname{PD}(\widetilde\alpha_i)$，定义
\[
  \operatorname{GW}^{\widetilde M}_{\widetilde A,k}
  (\widetilde\alpha_1,\ldots,\widetilde\alpha_k)
  :=\operatorname{GW}^{\widetilde M}_{\widetilde A,k}
  (\widetilde a_1,\ldots,\widetilde a_k).
\]

\begin{exercise}\label{ex:8-6-2}
证明：不变量 $\operatorname{GW}^{\widetilde M}_{\widetilde A,k}$ 与 $j_\Sigma$ 无关。

\emph{提示：}修改定理~\ref{thm:8-5-1} 的证明，使 $\Sigma$ 上的复结构也可以变化。
\end{exercise}

\begin{remark}\label{rem:8-6-3}
若 $\Sigma=S^2$，且 $(\widetilde M,\widetilde\omega)$ 本身\BookRecordIndex{T0058}半正，则对
$H_2(\widetilde M)$ 中所有\BookRecordIndex{T0285}球面类 $B$，定理~\ref{thm:7-1-1} 都定义了
$\widetilde M$ 的\BookRecordIndex{T0143}格罗莫夫--威滕不变量
$\operatorname{GW}^{\widetilde M}_{B,k}$。正如记号所示，当 $B$ 是\BookRecordIndex{T0189}截面类
$\widetilde A$ 时，这些先前定义的不变量与定理~\ref{thm:8-6-1} 的新不变量一致。在
这种情形，定理~\ref{thm:8-6-1} 的内容是说明：
$\operatorname{GW}^{\widetilde M}_{\widetilde A,k}$ 可以只用适配于\BookRecordIndex{T0279}纤维化的那一类
$\widetilde\omega$-\BookRecordIndex{T0372}驯服\BookRecordIndex{T0088}殆复结构，即形如 $\widetilde J_H$ 的\BookRecordIndex{T0088}殆复结构来计算。换言之，
这一类\BookRecordIndex{T0088}殆复结构足够丰富，能够保证\BookRecordIndex{T0166}横截性。由于投影
$\widetilde M\to S^2$ 是 $(\widetilde J_H,j_{S^2})$-全纯的，每条表示\BookRecordIndex{T0189}截面类的
$\widetilde J_H$-全纯曲线 $v:S^2\to\widetilde M$，投影到 $S^2$ 后都是次数为一的
全纯曲线。因此 $\operatorname{im}v$ 是一个截面。新定义已经内置了这一事实：\BookRecordIndex{T0246}模空间
$\mathcal M_{\Sigma,k}(\widetilde A;J,H)$ 只由带自然参数化的截面组成。

另一个特殊情形仍取 $\Sigma=S^2$，但\BookRecordIndex{T0279}纤维化是平凡的：
\[
  \widetilde M:=S^2\times M,
  \qquad \widetilde\omega=\operatorname{pr}^*\omega,
\]
其中 $\operatorname{pr}:\widetilde M\to M$ 是显然的投影。此时我们断言，定理
~\ref{thm:8-6-1} 的不变量与定理~\ref{thm:7-3-1} 定义的不变量
$\widetilde{\operatorname{GW}}^{M,I}_{A,k}$ 包含相同的信息。（不过要注意，现在的
观点稍有不同：此处通过 $\widetilde M^k$ 中的相交而非 $M^k$ 中的相交来计算不变量。）
为解释这一点，固定指标集 $I\subset\{1,\ldots,k\}$。我们断言
\begin{equation}\label{eq:8-6-2}
  \widetilde{\operatorname{GW}}^{M,I}_{A,k}
  (\alpha_1,\ldots,\alpha_k)
  =\operatorname{GW}^{\widetilde M}_{\widetilde A,k}
  (\widetilde\alpha_1,\ldots,\widetilde\alpha_k)
\end{equation}
对 $\alpha_i\in H_*(M)$ 成立，其中 $\widetilde\alpha_i\in H_*(\widetilde M)$ 定义为
\[
  \widetilde\alpha_i:=
  \begin{cases}
    \iota_*\alpha_i,&i\in I,\\
    [S^2]\times\alpha_i,&i\notin I.
  \end{cases}
\]
这里 $\iota:M\to\widetilde M$ 表示纤维的包含映射。对 Poincaré 对偶的上同调类而言，
这意味着，当 $i\in I$ 时
$\widetilde a_i=\operatorname{pr}^*a_i\smile\operatorname{PD}(\text{纤维})$，而当
$i\notin I$ 时 $\widetilde a_i=\operatorname{pr}^*a_i$。换言之，在 $M$ 的
$\widetilde{\operatorname{GW}}$ 不变量中固定由 $I$ 标号的\BookRecordIndex{T0066}标记点，对应于和位于某个
纤维中的循环 $\iota_*\alpha_i$ 相交。这个断言只是命题~\ref{prop:7-3-6} 的同调版本，
也可用同样的论证证明。注意，命题~\ref{prop:7-3-6} 中的\BookRecordIndex{T0058}半正性假设在这里已不再需要，
因为由定理~\ref{thm:8-6-1}，不变量
$\operatorname{GW}^{\widetilde M}_{\widetilde A,k}$ 已有良好定义。
\end{remark}

\begin{exercise}\label{ex:8-6-4}
由式~\eqref{eq:7-3-2}，式~\eqref{eq:8-6-2} 左边仅在
\[
  2nk-\sum_i\deg(\alpha_i)
  =2n+2c_1(A)+2(k-\#I)
\]
时才可能非零。验证这与右边非平凡所需的维数条件一致，并从几何上解释两个计数为何
对应。
\end{exercise}

现在说明一般\BookRecordIndex{T0156}哈密顿纤维化 $\widetilde M\to\Sigma$ 也有类似结论。仍用同调不变量来叙述
最为方便：
\[
  \operatorname{GW}^{\widetilde M}_{\widetilde A,k}:
  H_*(\widetilde M^k)\longrightarrow\mathbb Z.
\]
回忆习题~\ref{ex:8-5-4} 中定义的 $\widetilde M^k$ 的子流形
\[
  \widetilde M^k(\mathbf w)
  :=\bigl\{\widetilde{\mathbf x}
  =(\widetilde x_1,\ldots,\widetilde x_k)\in\widetilde M^k
  \bigm|\pi(\widetilde x_i)=w_i\text{ for }i\in I\bigr\}.
\]
我们将使用习题~\ref{ex:8-5-4} 构造的\BookRecordIndex{T0337}伪循环
\[
  \widetilde{\operatorname{ev}}_{\mathbf w}
  =\widetilde{\operatorname{ev}}_{\mathbf w,J,H}:
  \mathcal M_{\Sigma,k}(\widetilde A;\mathbf w,J,H)
  \longrightarrow\widetilde M^k(\mathbf w),
\]
并以
\[
  \iota_{\mathbf w}:\widetilde M^k(\mathbf w)\longrightarrow\widetilde M^k
\]
表示显然的包含映射。

\begin{proposition}\label{prop:8-6-5}
设 $(M,\omega)$、$(\widetilde M,\pi,\widetilde\omega)$、$\Sigma$ 和
$\widetilde A\in H_2(\widetilde M;\mathbb Z)$ 如定理~\ref{thm:8-6-1}。设 $k$ 为正整数，
$I\subset\{1,\ldots,k\}$，$\mathbf w=\{w_i\}_{i\in I}$ 是 $\Sigma$ 上两两不同的
点所成的元组，并且
$(J,H)\in\mathcal{JH}_{\mathrm{reg}}
(\widetilde M,\pi,\widetilde\omega;\mathbf w)$。则
\[
  \operatorname{GW}^{\widetilde M}_{\widetilde A,k}
  (\iota_{\mathbf w*}\widetilde\alpha)
  =\widetilde{\operatorname{ev}}_{\mathbf w}
    \mathbin{\boldsymbol\cdot}\widetilde\alpha
\]
对每个 $\widetilde\alpha\in H_*(\widetilde M^k(\mathbf w))$ 都成立。
\end{proposition}

\begin{proof}
习题~\ref{ex:8-5-4} 已经说明
$\widetilde{\operatorname{ev}}_{\mathbf w}$ 是\BookRecordIndex{T0337}伪循环，故右边有良好定义。不失一般性，
假设 $\widetilde\alpha$ 可由定向子流形
$X\subset\widetilde M^k(\mathbf w)$ 表示（见注~\ref{rem:6-5-3}）。由
引理~\ref{lem:6-5-5}，可以选取 $X$，使它与\BookRecordIndex{T0337}伪循环
$\widetilde{\operatorname{ev}}_{\mathbf w}$ \BookRecordIndex{T0281}强横截。于是，子流形
$\iota_{\mathbf w}(X)\subset\widetilde M^k$ 与\BookRecordIndex{T0337}伪循环
\[
  \widetilde{\operatorname{ev}}:
  \mathcal M_{\Sigma,k}(\widetilde A;J,H)\longrightarrow\widetilde M^k
\]
\BookRecordIndex{T0281}强横截。（这是因为 $\mathcal M_{\Sigma,k}(\widetilde A;J,H)$ 中每个元素都是截面；
因此，在 $w_i$ 处沿纤维法向方向所需的额外\BookRecordIndex{T0166}横截性，可以通过移动第 $i$ 个\BookRecordIndex{T0066}标记点得到。）
所以
$\iota_{\mathbf w}(X)\mathbin{\boldsymbol\cdot}\widetilde{\operatorname{ev}}
=X\mathbin{\boldsymbol\cdot}\widetilde{\operatorname{ev}}_{\mathbf w}$，命题得证。
\end{proof}

在解释这个结果的含义之前，先为 $\widetilde M$ 中通过计数截面得到的\BookRecordIndex{T0143}格罗莫夫--威滕
不变量引入如下记号。由于我们允许 $\pi_1(\Sigma)$ 在纤维的（上）同调上作非平凡作用，
这些不变量可能依赖于把各纤维等同于 $M$ 的方式。为简单起见（而且后文只会用到这一
情形），以下限制在 $I=\{1,\ldots,k\}$ 的情形。

\begin{definition}\label{def:8-6-6}
设 $(M,\omega)$、$(\widetilde M,\pi,\widetilde\omega)$、$\Sigma$ 和
$\widetilde A\in H_2(\widetilde M;\mathbb Z)$ 如定理~\ref{thm:8-6-1}。令 $k$ 为非负
整数，而 $\mathbf w=\{w_i\}_{1\le i\le k}$ 是 $\Sigma$ 上两两不同的点所成的 $k$ 元组。
对 $i=1,\ldots,k$，令 $\iota_i:M\to\widetilde M$ 为满足
$\iota_i(M)=\pi^{-1}(w_i)$ 的辛嵌入，并记
\[
  \iota:=\iota_1\times\cdots\times\iota_k:
  M^k\longrightarrow\widetilde M^k.
\]
定义如下群同态\sourcefootnote{原书把定义域印作单个 $H_*(M)$，但紧接着的定义以 $k$ 个同调类
$\alpha_1,\ldots,\alpha_k$ 为输入，严格说这是一个 $k$ 重多线性映射，等价地定义在
$H_*(M)^{\otimes k}$ 上；式中再由交叉积把纯张量送入 $H_*(M^k)$。}：
\[
  \operatorname{GW}^{\widetilde M,\mathbf w}_{\widetilde A,k}:
  H_*(M)^{\otimes k}\longrightarrow\mathbb Z
\] 其值为
\begin{equation}\label{eq:8-6-3}
  \operatorname{GW}^{\widetilde M,\mathbf w}_{\widetilde A,k}
  (\alpha_1,\ldots,\alpha_k)
  :=\operatorname{GW}^{\widetilde M}_{\widetilde A,k}
  \bigl(\iota_*(\alpha_1\times\cdots\times\alpha_k)\bigr).
\end{equation}
相应的上同调不变量通过 $M$ 中的 Poincaré 对偶定义：
\[
  \operatorname{GW}^{\widetilde M,\mathbf w}_{\widetilde A,k}
  (a_1,\ldots,a_k)
  :=\operatorname{GW}^{\widetilde M,\mathbf w}_{\widetilde A,k}
  \bigl(\operatorname{PD}_M(a_1),\ldots,\operatorname{PD}_M(a_k)\bigr),
  \qquad a_1,\ldots,a_k\in H^*(M).
\]
\end{definition}

命题~\ref{prop:8-6-5} 表明，式~\eqref{eq:8-6-3} 的不变量既可以表示为
$\widetilde{\operatorname{ev}}$ 与 $\widetilde M^k$ 中类
$\iota_*(\alpha_1\times\cdots\times\alpha_k)$ 的一个代表元之交数，也可以表示为
$\iota^{-1}\circ\widetilde{\operatorname{ev}}_{\mathbf w}$ 与 $M^k$ 中类
$\alpha_1\times\cdots\times\alpha_k$ 的一个代表元之交数。其作用在于：当
\[
  (J,H)\in\mathcal{JH}_{\mathrm{reg}}
  (\widetilde M,\pi,\widetilde\omega;\mathbf w)
\]
且 $\widetilde X_i\subset\pi^{-1}(w_i)$ 是表示同调类 $\iota_{i*}\alpha_i$ 的子流形时，
可以在各自的纤维中（而不是在环境流形 $\widetilde M$ 中）扰动
$\widetilde X_i$，使乘积
\[
  \widetilde X:=\widetilde X_1\times\cdots\times\widetilde X_k
\]
与\BookRecordIndex{T0337}伪循环 $\widetilde{\operatorname{ev}}$ \BookRecordIndex{T0281}强横截。因此，这些不变量具有与
$\widetilde{\operatorname{GW}}^{M,I}_{A,k}$ 相同的形式性质；特别地，它们满足
（分裂）公理的一个类似版本，见定理~11.4.6。这是讨论 \BookRecordIndex{T0047}Seidel 表示时的一个重要要素。

\paragraph{例子。}
现在给出一些例子。前两个本质上用到了\BookRecordIndex{T0156}哈密顿纤维化，最后一个则没有。最后一个例子
利用了如下事实：当复结构 $j_\Sigma$ 固定时，我们现在已经精确定义了高亏格不变量；
它将展示这种不变量的一些新特征。

\begin{example}[\BookRecordIndex{T0032}Lefschetz 数]\label{ex:8-6-7}
这个例子出自 Ionel--Parker \srcite{199}。它与辛\BookRecordIndex{T0132}弗洛尔同调、三维流形的
Seiberg--Witten 不变量以及 Turaev 版本的 Reidemeister \BookRecordIndex{T0248}挠率有一些有趣联系。设
$(M,\omega)$ 是满足强\BookRecordIndex{T0058}半正性条件~\eqref{eq:8-5-1} 的紧\BookRecordIndex{T0358}辛流形，并设
\[
  f:M\longrightarrow M
\]
为辛同胚。把
\[
  M_f:=\mathbb R\times M\big/\bigl((t,f(x))\sim(t+1,x)\bigr)
\]
称为\BookRecordIndex{T0381}映射圆。考虑\BookRecordIndex{T0200}局部哈密顿纤维化
\[
  \widetilde M_f:=S^1\times M_f\longrightarrow\mathbb T^2.
\]
由于 $f$ 是辛同胚，$\omega$ 在投影 $S^1\times\mathbb R\times M\to M$ 下的拉回
下降为\BookRecordIndex{T0224}联络形式 $\widetilde\omega_f\in\Omega^2(\widetilde M_f)$。现在任取
$M$ 上的 $\omega$-\BookRecordIndex{T0372}驯服\BookRecordIndex{T0088}殆复结构族
$J=\{J_t\}_{t\in\mathbb R}$，满足 $J_{t+1}=f^*J_t$。则 $J$ 是
$\widetilde M_f$ 上的\BookRecordIndex{T0306}竖直殆复结构，因而与 $\widetilde\omega_f$ 一起确定
$\widetilde M_f$ 上的\BookRecordIndex{T0088}殆复结构 $\widetilde J_f$。每个不动点
$x\in\operatorname{Fix}(f)$ 都确定一个截面
$\widetilde u_x(s,t):=x$，从而确定一个同调类
\[
  \widetilde A_x:=[\widetilde u_x]\in H_2(\widetilde M_f;\mathbb Z).
\]
由定义容易看出，竖直 Fredholm 算子 $D_{\widetilde u_x}$ 的指标为零。这个算子双射
当且仅当 $x$ 是非退化不动点，即
\[
  \det\bigl(\mathbf 1-df(x)\bigr)\ne0.
\]
此外，可以证明，截面 $\widetilde u_x$ 对不变量的贡献等于不动点指数
$\operatorname{sign}\det(\mathbf 1-df(x))$。由于 $\widetilde\omega_f$ 的曲率为零，
引理~\ref{lem:8-2-9} 的\BookRecordIndex{T0307}竖直能量恒等式说明，$\widetilde M_f$ 中每个表示类
$\widetilde A_x$ 的 $\widetilde J_f$-全纯截面都必须是水平的，因而对应于 $f$ 的一个
不动点。

现在假设 $f$ 的所有不动点都非退化，并设
$\widetilde A\in H_2(\widetilde M_f;\mathbb Z)$ 是满足
\[
  \pi_*\widetilde A=[\Sigma],\qquad
  \widetilde\omega_f(\widetilde A)=0,\qquad
  c_1^{\mathrm{Vert}}(\widetilde A)=0
\]
的同调类。称这样的类为\emph{\BookRecordIndex{T0316}水平类}。上述讨论表明，\BookRecordIndex{T0316}水平类 $\widetilde A$ 的无标记
\BookRecordIndex{T0143}格罗莫夫--威滕不变量等于 \BookRecordIndex{T0032}Lefschetz 数：
\[
  \operatorname{GW}^{\widetilde M_f}_{\widetilde A,0}(1)
  =L(f;\widetilde A)
  :=\sum_{\substack{x\in\operatorname{Fix}(f)\\
                    \widetilde A_x=\widetilde A}}
    \operatorname{sign}\det\bigl(\mathbf 1-df(x)\bigr).
\]
对所有\BookRecordIndex{T0316}水平类 $\widetilde A$ 求和，便得到完整的 \BookRecordIndex{T0032}Lefschetz 数
\[
  \sum_{\widetilde A\text{ 为\BookRecordIndex{T0316}水平类}}
  \operatorname{GW}^{\widetilde M_f}_{\widetilde A,0}(1)
  =L(f)=\operatorname{trace}(f_{\mathrm{ev}})
       -\operatorname{trace}(f_{\mathrm{odd}}),
\]
其中
$f_{\mathrm{odd}}:H_{\mathrm{odd}}(M)\to H_{\mathrm{odd}}(M)$ 表示 $f$ 在奇数维同调上
诱导的同态，$f_{\mathrm{ev}}$ 的定义类似。$f$ 的 zeta 函数与迭代 $f^k$ 的
\BookRecordIndex{T0032}Lefschetz 数之间有关系
\[
  \zeta_f(t)
  =\exp\left(\sum_{k=1}^{\infty}\frac{L(f^k)t^k}{k}\right)
  =\frac{\det(\mathbf 1-tf_{\mathrm{odd}})}
         {\det(\mathbf 1-tf_{\mathrm{ev}})}.
\]
若要用 $\widetilde M_f$ 的\BookRecordIndex{T0143}格罗莫夫--威滕不变量的生成函数来解释这个 zeta 函数，
就必须考虑同调类
$\widetilde A\in H_2(\widetilde M_f;\mathbb Z)$
\sourcefootnote{原书此处把同调群印作 $H_2(\widetilde M;\mathbb Z)$；但本例的总空间
从头至尾都是 $\widetilde M_f$，且紧接着也在其中考察全纯曲线，故补回下标 $f$。}，满足
\[
  \pi_*\widetilde A=k[\Sigma]
\] 并考察 $\widetilde M_f$ 中截面以外的全纯曲线。
Ionel 和 Parker 在 \srcite{199} 中完成了这项工作；本书不讨论这一推广。
\end{example}

\begin{example}[\BookRecordIndex{T0047}Seidel 表示]\label{ex:8-6-8}
设 $(M,\omega)$ 是紧单调\BookRecordIndex{T0358}辛流形。Seidel 在 \srcite{363} 中为哈密顿辛同胚的每个环路
\[
  \psi=\{\psi_t\}_{t\in\mathbb R/\mathbb Z}
\]
配上一个称为 \emph{\BookRecordIndex{T0046}Seidel 元}的上同调类
$\mathcal S(\psi)\in H^*(M)$。每个这样的环路都像式~\eqref{eq:8-2-6} 那样定义一个
\BookRecordIndex{T0156}哈密顿纤维化 $\widetilde M_\psi$，而类 $\mathcal S(\psi)$ 是通过让
$\widetilde M_\psi$ 的全纯截面\BookRecordIndex{T0246}模空间与一个纤维相交得到的。更准确地说，定义
\[
  \mathcal S(\psi):=\sum_{\widetilde A}\mathcal S_{\widetilde A}(\psi),
\]
其中上同调类
$\mathcal S_{\widetilde A}(\psi)\in H^{2c}(M)$ 与如下\BookRecordIndex{T0337}伪循环在 $M$ 中表示的同调类
Poincaré 对偶：
\[
  \widetilde{\operatorname{ev}}_{\mathbf w,J,H}:
  \mathcal M(\widetilde A;J,H)\longrightarrow M.
\]
这里\BookRecordIndex{T0246}模空间由类 $\widetilde A$ 中带一个\BookRecordIndex{T0150}固定标记点 $\mathbf w$ 的截面组成；
$-c:=c_1^{\mathrm{Vert}}(\widetilde A)$ 是 $\widetilde A$ 的竖直 Chern 数，并且我们把
$\widetilde M_\psi$ 在 $\mathbf w$ 上的纤维等同于 $M$。这个和遍历所有满足
\[
  \widetilde A\in H_2(\widetilde M_\psi;\mathbb Z),
  \qquad \pi_*\widetilde A=[S^2]
\]
的同调类。故 $\widetilde A$ 只有在 $-2n\le-c\le0$ 时才有贡献；$(M,\omega)$ 的
单调性说明这个和是有限的。等价地，可以用定义~\ref{def:8-6-6} 的\BookRecordIndex{T0143}格罗莫夫--威滕
不变量
\[
  \operatorname{GW}^{\widetilde M_\psi,\mathbf w}_{\widetilde A,1}:
  H^{2n-2c}(M)\longrightarrow\mathbb Z
\]
按下式表示 $\mathcal S_{\widetilde A}(\psi)$：
\begin{equation}\label{eq:8-6-4}
  \mathcal S_{\widetilde A}(\psi)
  :=\sum_{\nu,\mu}
  \operatorname{GW}^{\widetilde M_\psi,\mathbf w}_{\widetilde A,1}(e_\nu)
  g^{\nu\mu}e_\mu
  \in H^{2c}(M).
\end{equation}
这里仍把 $M$ 与 $\mathbf w$ 上的纤维等同。照例，
$e_0,\ldots,e_m$ 表示 $H^*(M)$ 的一组基，而 $g^{\nu\mu}$ 是矩阵
$g_{\nu\mu}$ 的逆矩阵，其中
\[
  g_{\nu\mu}:=\int_M e_\nu\smile e_\mu.
\]

对一般流形 $(M,\omega)$，必须按不同的类 $\widetilde A$ 分开
$\mathcal S(\psi)$ 的各项贡献。最自然的做法是使用\BookRecordIndex{T0229}量子上同调。事实上，第~11.4 节将
说明，$\mathcal S(\psi)$ 可以解释为 $(M,\omega)$ 的\BookRecordIndex{T0229}量子上同调环
$QH^*(M)$ 中的可逆元。此外，映射
\[
  \pi_1(\operatorname{Ham}(M,\omega))\longrightarrow H^*(M),
  \qquad \psi\longmapsto\mathcal S(\psi)
\]
是从 $(M,\omega)$ 的哈密顿辛同胚群之基本群到 $QH^*(M)$ 可逆元乘法群的群同态。让
$\mathcal S(\phi)$ 通过量子乘法作用在 $\Lambda$-模 $QH^*(M)$ 上，就得到
$\pi_1(\operatorname{Ham}(M,\omega))$ 在 $QH^*(M)$ 作为左 $QH^*(M)$-模的自同构群中
的一个表示，称为 \emph{\BookRecordIndex{T0047}Seidel 表示}。

要使上述陈述成立，当环路
$\psi=\{\psi_t\}_{t\in\mathbb R/\mathbb Z}$ 可缩时，就应有
$\mathcal S(\psi)=1$。为验证这一点，注意此时\BookRecordIndex{T0156}哈密顿纤维化
$\widetilde M_\psi$ 逐纤维辛同胚于乘积 $S^2\times M$。因此，由习题~\ref{ex:7-3-4}，
当 $\widetilde A\ne0$ 时不变量为零；而由习题~\ref{ex:7-3-5}，
\[
  \operatorname{GW}^{\widetilde M_\psi,\mathbf w}_{0,1}(a)=\int_Ma.
\]
所以 $\mathcal S(\psi)=\mathcal S_0(\psi)=1$。第~\ref{sec:9-6} 节还会回到这个例子。
\end{example}

现在考虑 $[\omega]$ 在 $\pi_2(M)$ 上消失这一特殊情形。此时对任何 $J$ 都不存在
$J$-全纯球面，因而 $M$ 的\BookRecordIndex{T0229}量子上同调环结构就是通常的杯积。在这种情形，\BookRecordIndex{T0047}Seidel 表示的
乘法性质具有如下形式。

\begin{theorem}[Seidel]\label{thm:8-6-9}
设 $(M,\omega)$ 是紧\BookRecordIndex{T0358}辛流形，并且 $[\omega]$ 在 $\pi_2(M)$ 上消失。则对
$M$ 的哈密顿辛同胚的任意两个环路
$\psi=\{\psi_t\}_{t\in\mathbb R/\mathbb Z}$ 和
$\phi=\{\phi_t\}_{t\in\mathbb R/\mathbb Z}$，都有
\[
  \mathcal S(\psi\phi)=\mathcal S(\psi)\smile\mathcal S(\phi).
\]
\end{theorem}

\begin{proof}
见第~11.4 节。
\end{proof}

下面这个有趣的推论将在第~\ref{sec:9-1} 节得到更完整的解释。当第一 Chern 类在 $\pi_2(M)$ 上
消失时，Schwarz 在 \srcite{360} 中利用可追溯至 Seidel 的思想证明了它。主要步骤是
说明，对这类流形，\BookRecordIndex{T0046}Seidel 元 $\mathcal S(\psi)$ 的首项
$\mathcal S(\psi)_0\in H^0(M)$ 等于 $\pm1$。

\begin{corollary}[Schwarz]\label{cor:8-6-10}
设 $(M,\omega)$ 是紧\BookRecordIndex{T0358}辛流形，且 $[\omega]$ 在 $\pi_2(M)$ 上消失；设
$\{\psi_t\}_{t\in\mathbb R/\mathbb Z}$ 是 $M$ 的哈密顿辛同胚环路。令
$H_t:M\to\mathbb R$ 是生成 $\psi_t$、均值为零的哈密顿量，即
\[
  \partial_t\psi_t=X_{H_t}\circ\psi_t,\qquad
  \iota(X_{H_t})\omega=dH_t,\qquad
  \int_M H_t\omega^n=0.
\]
设 $\mathbb D\subset\mathbb C$ 是闭单位圆盘，而 $u:\mathbb D\to M$ 是满足
\[
  u(e^{2\pi it})=\psi_t(x_0),\qquad x_0:=u(1)
\]
的光滑映射。则
\[
  \int_{\mathbb D}u^*\omega
  +\int_0^1H_t\bigl(\psi_t(x_0)\bigr)\,dt=0.
\]
\end{corollary}

\begin{proof}
设 $\pi:\widetilde M_\psi\to S^2$ 是式~\eqref{eq:8-2-6} 的\BookRecordIndex{T0156}哈密顿纤维化，并令
$\widetilde\omega\in\Omega^2(\widetilde M_\psi)$ 是注~\ref{rem:8-2-11} 中的\BookRecordIndex{T0224}联络形式，
其数据为
\[
  F^\pm=0,\qquad
  G^+_{s,t}=-\rho(s)H_t,\qquad
  G^-_{s,t}=-\rho(s)H_{-t}\circ\psi_{-t},
\]
其中 $\rho:\mathbb R\to[0,1]$ 是光滑\BookRecordIndex{T0188}截断函数，满足：当 $s\le-1$ 时
$\rho(s)=0$，当 $s\ge1$ 时 $\rho(s)=1$，并且对所有 $s$ 都有
$\rho(s)+\rho(-s)=1$。不失一般性，假设
\[
  u(e^{2\pi(s+it)})=\psi_t(x_0)\qquad(s\ge-1),
\]
并以
\[
  \widetilde u^+(s,t):=u(e^{2\pi(s+it)}),\qquad
  \widetilde u^-(s,t):=
  \psi_{-t}\bigl(u(e^{2\pi(-s-it)})\bigr)
\]
定义截面 $\widetilde u:S^2\to\widetilde M_\psi$。则
\[
  a(\psi):=-\int_{S^2}\widetilde u^*\widetilde\omega
  =-\int_{\mathbb D}u^*\omega
   -\int_0^1H_t\bigl(\psi_t(x_0)\bigr)\,dt.
\]
由于 $[\omega]$ 在 $\pi_2(M)$ 上消失，这个积分与 $\widetilde M_\psi$ 的截面
$\widetilde u$ 无关，因而尤其与点 $x_0$ 的选择无关。我们必须证明 $a(\psi)=0$。

固定 $\widetilde M_\psi$ 上的\BookRecordIndex{T0306}竖直殆复结构 $J$。给定任意水平 1-形式
$\widetilde\sigma_K\in\Omega^1(\widetilde M_\psi)$
\sourcefootnote{原书此处印作
$\widetilde\sigma_K\in H^1(\widetilde M_\psi)$；但紧接着要取
$d\widetilde\sigma_K$ 并用它改变\BookRecordIndex{T0224}联络形式，故该对象唯一应属于
$\Omega^1(\widetilde M_\psi)$。}，以 $\widetilde J_K$ 表示由 $J$ 与
$\widetilde\omega-d\widetilde\sigma_K$ 在 $\widetilde M_\psi$ 上诱导的\BookRecordIndex{T0088}殆复结构，并以
$\mathcal M(J,K)$ 表示 $\widetilde M_\psi$ 的所有 $\widetilde J_K$-全纯截面所成的
\BookRecordIndex{T0246}模空间。因为 $[\omega]$ 在 $\pi_2(M)$ 上消失，这个\BookRecordIndex{T0246}模空间是紧的。若 $K$ 在
式~\eqref{eq:8-3-2} 的意义下正则，则由定理~\ref{thm:8-3-1}，
$\mathcal M(J,K)$ 是紧定向光滑流形。它的不同连通分支可能具有不同维数：与同调类
$\widetilde A\in H_2(\widetilde M_\psi;\mathbb Z)$ 对应的连通分支维数为
\[
  2n+2c_1^{\mathrm{Vert}}(\widetilde A).
\]
记 $\mathcal M_0(J,K)$ 为 Chern 数为零的连通分支之并。在正则情形，\BookRecordIndex{T0283}求值映射
\[
  \operatorname{ev}_K:\mathcal M(J,K)\longrightarrow M
\]
表示同调类 $\operatorname{PD}(\mathcal S(\psi))$，而限制后的\BookRecordIndex{T0283}求值映射
\[
  \operatorname{ev}_{K,0}:\mathcal M_0(J,K)\longrightarrow M
\]
表示首项
\[
  \mathcal S(\psi)_0
  =\deg(\operatorname{ev}_{K,0})\,1\in H^0(M).
\]
由定理~\ref{thm:8-6-9}，
$\mathcal S(\psi)\smile\mathcal S(\psi^{-1})=1$。这只有在\BookRecordIndex{T0283}求值映射
$\operatorname{ev}_{K,0}:\mathcal M_0(J,K)\to M$ 的度数为 $+1$ 或 $-1$ 时才可能。

现在考虑\BookRecordIndex{T0283}求值映射
\[
  \widetilde{\operatorname{ev}}_{K,0}:
  \mathcal M_0(J,K)\times S^2\longrightarrow\widetilde M_\psi,
\]
并对 $\widetilde\omega^{n+1}\in
\Omega^{2n+2}(\widetilde M_\psi)$ 的拉回积分，得到
\[
\begin{aligned}
  a(\psi)\deg(\operatorname{ev}_{K,0})
  \int_M\frac{\omega^n}{n!}
  &=
  a(\psi)\int_{\mathcal M_0(J,K)}
  \frac{\operatorname{ev}_{K,0}^*\omega^n}{n!}\\
  &=-\int_{\mathcal M_0(J,K)\times S^2}
  \frac{\widetilde{\operatorname{ev}}_{K,0}^*
  \widetilde\omega^{n+1}}{(n+1)!}\\
  &=-\deg(\widetilde{\operatorname{ev}}_{K,0})
  \int_{\widetilde M_\psi}
  \frac{\widetilde\omega^{n+1}}{(n+1)!}\\
  &=0.
\end{aligned}
\]
第一个和第三个等式来自度数定理，第二个等式来自下面的习题~\ref{ex:8-6-11}，最后一个
来自习题~\ref{ex:8-1-3}。因为
$\deg(\operatorname{ev}_{K,0})\ne0$，所以得到所需的 $a(\psi)=0$。推论得证。
\end{proof}

\begin{exercise}\label{ex:8-6-11}
设 $X$ 是紧定向 $2n$ 维流形，而
$\tau\in\Omega^2(X\times S^2)$ 是闭 2-形式。设
$\iota_X:X\to X\times S^2$ 与
$\iota_{S^2}:S^2\to X\times S^2$ 是嵌入
\[
  \iota_X(x):=(x,z_0),\qquad
  \iota_{S^2}(z):=(x_0,z).
\]
证明
\[
  \int_{X\times S^2}\frac{\tau^{n+1}}{(n+1)!}
  =
  \left(\int_X\frac{\iota_X^*\tau^n}{n!}\right)
  \left(\int_{S^2}\iota_{S^2}^*\tau\right).
\]

\emph{提示：}由于 $S^2$ 单连通，Künneth 公式断言，存在闭 2-形式
$\tau_X\in\Omega^2(X)$ 和 $\tau_{S^2}\in\Omega^2(S^2)$，使
\[
  \tau-\pi_X^*\tau_X-\pi_{S^2}^*\tau_{S^2}
\]
为\BookRecordIndex{T0403}正合形式。
\end{exercise}

\begin{example}\label{ex:8-6-12}
考虑 $M=\mathbb{CP}^n$ 的情形，并设 $\Sigma$ 是亏格为 $g$ 的黎曼曲面。取
$d\ge0$，使
\[
  m:=(n+1)(d+1-g)+g-1\ge0.
\]
这个数是次数为 $d$ 的全纯映射 $u:\Sigma\to\mathbb{CP}^n$ 之\BookRecordIndex{T0246}模空间虚维数的一半。
令 $c\in H^2(\mathbb{CP}^n;\mathbb Z)$ 表示正生成元，并选取整数
$m_1,\ldots,m_k\in\mathbb N$，满足
\[
  0\le m_i\le n,\qquad \sum_{i=1}^k m_i=m.
\]
则
\begin{equation}\label{eq:8-6-5}
  \operatorname{GW}^{\mathbb{CP}^n,\{1,\ldots,k\}}_{\Sigma,dL,k}
  (c^{m_1},\ldots,c^{m_k})=(n+1)^g.
\end{equation}
正如记号所示，这个不变量可以只用到 $\mathbb{CP}^n$ 的亏格 $g$ \BookRecordIndex{T0340}稳定映射来定义。
不过，把这些\BookRecordIndex{T0340}稳定映射等同于平凡丛 $\Sigma\times\mathbb{CP}^n$ 的截面有时很有用。
（式~\eqref{eq:8-6-5} 的证明见 Bertram--Daskalopoulos--Wentworth \srcite{36}，
另见 Cieliebak--Gaio--Mundet--Salamon \srcite{67}。）

更仔细地考察 $n=g=d=1$ 这一特殊情形很有意思。此时 $m=2$，可以取
$k=2$ 且 $m_1=m_2=1$。由式~\eqref{eq:8-6-5}，这个不变量等于 $2$，尽管不存在
从二维环面到二维球面的次数为一的全纯映射。不过，次数 $d=1$、亏格 $g=1$ 的\BookRecordIndex{T0340}稳定映射
确实构成一个非空\BookRecordIndex{T0246}模空间：$\mathbb T^2$ 上的复结构固定，且有 $k=2$ 个\BookRecordIndex{T0066}标记点
$z_1,z_2$，它们到环面的投影分别固定在 $w_1,w_2\in\mathbb T^2$。每个这样的\BookRecordIndex{T0340}稳定映射
都由一个亏格一的常值分支和一条次数为一的非常值\BookRecordIndex{T0386}有理曲线组成；后者在环面上的某一点
附着于前者，我们把这个点称为\BookRecordIndex{T0191}节点。若\BookRecordIndex{T0191}节点属于 $\{w_1,w_2\}$，还可能有一个包含相应
\BookRecordIndex{T0066}标记点的\BookRecordIndex{T0384}幽灵气泡。不过要注意，此时两个\BookRecordIndex{T0066}标记点在 $M=S^2$ 中的像相同。在乘积
$\mathbb T^2\times S^2$ 中，常值分支表现为常值映射的图像，而\BookRecordIndex{T0386}有理曲线表现为\BookRecordIndex{T0191}节点上方纤维中
的全纯球面。

这里共有五个\BookRecordIndex{T0072}层：一个\BookRecordIndex{T0072}层中，环面上的\BookRecordIndex{T0191}节点不属于 $\{w_1,w_2\}$；其余四个\BookRecordIndex{T0072}层中，\BookRecordIndex{T0191}节点为
$w_1$ 或 $w_2$，而相应的\BookRecordIndex{T0066}标记点位于\BookRecordIndex{T0386}有理曲线上或\BookRecordIndex{T0384}幽灵气泡上。由此容易看出，在这个
空间中，恰有两个\BookRecordIndex{T0340}稳定映射把两个\BookRecordIndex{T0066}标记点送到预先指定的两个不同点
$x_1,x_2\in S^2$。其中一个把环面映到 $x_1$，\BookRecordIndex{T0191}节点为 $w_2$，第二个\BookRecordIndex{T0066}标记点位于球面上
并映到 $x_2$；另一个把环面映到 $x_2$，\BookRecordIndex{T0191}节点为 $w_1$，第一个\BookRecordIndex{T0066}标记点位于球面上并映到
$x_1$。一般的\BookRecordIndex{T0157}哈密顿扰动会给出扰动\BookRecordIndex{T0209}柯西--黎曼方程的两个正则解
\[
  u:\mathbb T^2\longrightarrow S^2,\qquad
  u(w_1)=x_1,\quad u(w_2)=x_2.
\]

计数 $\mathbb{CP}^2$ 中的全纯曲线时也有类似情形。经典枚举不变量给出通过适当数目
点的亏格 $g$、次数 $d$ 曲线之数目。它们与计数通过这些约束的所有\BookRecordIndex{T0340}稳定映射之相应
\BookRecordIndex{T0143}格罗莫夫--威滕不变量有关，但除亏格零的情形外，两者不必相同。造成差异的一个原因，
是上例中的幽灵（即常值）环面所造成的影响。进一步讨论可见 Ionel \srcite{196} 和
Zinger \srcite{426,427} 的论文。
\end{example}

\begin{remark}\label{rem:8-6-13}
上述\BookRecordIndex{T0156}哈密顿纤维化的讨论，大多集中于能阐明纤维 $M$ 性质的那些理论侧面。不过，以任意
光滑流形为底的\BookRecordIndex{T0156}哈密顿纤维化也有一套有趣的一般理论。这里的一个公开问题是
Lalonde--McDuff \srcite{227} 提出的 $c$-分裂猜想：一个光滑\BookRecordIndex{T0156}哈密顿纤维化的（紧）总空间
之有理上同调，作为加法群，是否同构于底空间与纤维的上同调之张量积？（注
~\ref{rem:8-4-4} 已经提到，当底空间为 $S^2$ 时，可以用 \BookRecordIndex{T0047}Seidel 表示证明这一点；
凯勒 情形也成立。）

更有意思的是辛 \BookRecordIndex{T0033}Lefschetz 纤维化的研究。这类\BookRecordIndex{T0279}纤维化是光滑映射
\[
  f:(M,\omega)\longrightarrow S^2,
\]
其纤维除有限个孤立奇点外都是辛的；在每个奇点附近，局部模型为
\[
  (z_1,\ldots,z_n)\longmapsto\sum_{i=1}^n z_i^2,
\]
其中 $(z_1,\ldots,z_n)\in\mathbb C^n$ 是局部 \BookRecordIndex{T0081}达布坐标卡中的坐标。因此，它们是
Morse 函数的复版本；Donaldson \srcite{85} 证明，只要
$[\omega]\in H^2(M;\mathbb Q)$，任意闭\BookRecordIndex{T0358}辛流形经适当爆破后都存在这种\BookRecordIndex{T0279}纤维化。辛
\BookRecordIndex{T0033}Lefschetz 纤维化在四维最为有用，因为此时纤维是二维的。例如，
Auroux--Muñoz--Presas \srcite{29} 已用它们研究辛四维流形中拉格朗日球面的同痕类以及
\BookRecordIndex{T0365}辛映射类群。
\end{remark}
